
\exercise{Codimension example \label{exManifoldExample}}{2}
Consider the system
\eqn{
\dot{x} = f(x,a,b) = a - bx^2. 
}
For $b>0$ find the steady states and their bifurcation points. Plot the bifurcation points in a 2-parameter bifurcation diagram, where the axes are the parameters $a$ and $b$. 

\solution
We compute the steady states
\eqa{
0 &=& a-bx^2 \\
 bx^2 &=& a\\
x^2 &=& a/b \\
x &=& \pm \sqrt{a/b}
}
So for $b>0$ and $a>0$ we have two branches of steady states $x_{1,2}=\pm\sqrt{a/b}$. At $a=0$ these two branches collide in a square root 
shape and annihilate each other so that there are no steady states left. 
This already shows that there is a fold bifurcation at $a=0$ regardless of $b$, 
as long as $b>0$. If we want to do another test we could compute 
\eq{
f_{\rm x} = \frac{\partial \dot{x}}{\partial x} = -2xb = \pm 2\sqrt{ab}
}
which becomes zero at $b=0$, confirming the presence of the bifurcation point.
We can thus plot the two-parameter diagram from the text:

FIGURE

As a bonus it is interesting to explore what happens at $b=0$. The line $b=0$ is also a bifurcation line as $f_{\rm x}=0$. However it is a very weird bifurcation as our steady states disappear to infinity. For $b<0$ the steady states reappear, but thy now exist for $a<0$. So this bifurcation line is none of the bifurcations that we studied so far but a very highly degenerate case, where almost all derivatives of $f$ become zero. such highly degenerate bifurcations can appear in very simple examples, because the most simple systems tend to be very special, and hence degenerate.

The point $a=b=0$ where the degenerate bifurcation intersects the fold is an example of a (highly degenerate) codimension-2 bifurcation where the fold bifurcation changes its nature from creating two steady states on the positive-$a$ side, to creating two steady states on the negative-$a$ side of the bifurcation manifold. 
